\section{PPT maps and the represented class}\label{sec:foundations}

All matrix algebras in the constructions are finite-dimensional over
$\C$, with fixed orthonormal bases and the induced product bases.  We write
$X^*$ for conjugate transpose, $\overline X$ for entrywise conjugation, and
$T_d(X)=X^{\T}$ for the complex-linear transpose on $\Mat d$.  For a linear
map $F$, its Hilbert--Schmidt adjoint $F^\dagger$ is characterized by
$\tr(Y^*F(X))=\tr((F^\dagger(Y))^*X)$.  If $V$ is a matrix, then
$\Ad_V(X)=VXV^*$.

For $F:\Mat a\to\Mat b$ our Choi convention is
\begin{equation}\label{eq:choi-definition}
 J(F)=\sum_{i,j=1}^a E_{ij}\otimes F(E_{ij})
     =(\mathrm{id}_a\otimes F)(\Omega_a\Omega_a^*),
 \qquad \Omega_a=\sum_{i=1}^a e_i\otimes e_i.
\end{equation}
In particular, $\norm{\Omega_a}^2=a$; the vector is not normalized.  The
first tensor factor is the input reference.  This convention fixes every
dimension factor in the transfer and channel formulas below.
It is the map--operator correspondence associated with Jamiołkowski and
Choi \cite{Jamiolkowski1972,Choi1975}.

\begin{lemma}\label{lem:ppt-stability}
Let $F:\Mat a\to\Mat b$ be CP.  Then $F^\dagger$ and
$T_b\circ F\circ T_a$ are CP.  Moreover,
\[
 F\circ T_a\text{ is CP}\quad\Longleftrightarrow\quad
 T_b\circ F\text{ is CP}.
\]
If $F$ is PPT, so are $T_b\circ F$, $F\circ T_a$, its adjoint, its
simultaneous transpose $T_b\circ F\circ T_a$, and its transpose-adjoint
$F^\sharp:=T_a\circ F^\dagger\circ T_b$.  Tensor products of PPT maps
are PPT.  If $P$ and $Q$ are compatible CP maps, then $Q\circ F\circ P$
is PPT whenever $F$ is PPT.
\end{lemma}

\begin{proof}
Choose Kraus matrices $A_r\in\C^{b\times a}$ such that
$F(X)=\sum_r A_rXA_r^*$.  Then
\[
 F^\dagger(Y)=\sum_r A_r^*YA_r,
 \qquad
 (T_b\circ F\circ T_a)(X)
       =\sum_r\overline{A_r}X\overline{A_r}^{\,*},
\]
which proves both CP assertions.  Applying the simultaneous-transpose
assertion to $F\circ T_a$ or to $T_b\circ F$ proves their equivalence.
This argument works with rectangular Kraus matrices.
The operator-sum representation is standard \cite{Kraus1971,Choi1975}.

Transpose is self-adjoint for the Hilbert--Schmidt inner product.  Hence,
when $F$ is PPT,
\[
 T_a\circ F^\dagger=(F\circ T_a)^\dagger
 \quad\text{and}\quad
 F^\dagger\circ T_b=(T_b\circ F)^\dagger
\]
are CP.  This proves the PPT assertion for the adjoint, and simultaneous
transposition gives the simultaneous transpose and transpose-adjoint.
The output transpose of $T_b\circ F$ is $F$, while that of
$F\circ T_a$ is $T_b\circ F\circ T_a$; all these maps are CP.  This proves
the remaining PPT assertions.  With product bases,
$T_{b_1b_2}=T_{b_1}\otimes T_{b_2}$, so the tensor-product assertion follows
from closure of CP maps under tensor products.  Finally, for
$P:\Mat{a'}\to\Mat a$ and $Q:\Mat b\to\Mat{b'}$,
\[
 T_{b'}\circ Q\circ F\circ P
 = (T_{b'}\circ Q\circ T_b)\circ(T_b\circ F)\circ P
\]
is a composition of CP maps.  The untransposed composition is CP as well.
\end{proof}

The input- and output-transpose formulations of PPT will therefore be used
interchangeably.  No trace condition is included in this terminology.
The elementary closure facts above are also recorded, in related forms, in
\cite[Lemma 2.1]{KMP2018}.

\begin{lemma}\label{lem:choi-separable}
A linear map $F:\Mat a\to\Mat b$ is CP if and only if $J(F)\succeq0$.
For a CP map, the following are equivalent:
\begin{enumerate}
 \item $J(F)$ is separable;
 \item for every $k\geq1$ and every $W\succeq0$ in $\Mat k\otimes\Mat a$,
       $(\mathrm{id}_k\otimes F)(W)$ is separable.
\end{enumerate}
Also,
\[
 J(T_b\circ F)=(\mathrm{id}_a\otimes T_b)J(F),\qquad
 J(F\circ T_a)=(T_a\otimes\mathrm{id}_b)J(F).
\]
\end{lemma}

\begin{proof}
The forward CP implication follows from \eqref{eq:choi-definition}.
Conversely, write a positive Choi matrix as
$J(F)=\sum_r v_rv_r^*$.  There is a unique matrix
$A_r\in\C^{b\times a}$ with $v_r=(I_a\otimes A_r)\Omega_a$.
Comparison of the $(i,j)$ blocks gives
$F(E_{ij})=\sum_r A_rE_{ij}A_r^*$ and hence the Kraus representation of
$F$.  This is Choi's finite-dimensional criterion \cite{Choi1975}.

If $J(F)$ is separable, spectral decompositions of its positive factors
write it as $\sum_r a_ra_r^*\otimes b_rb_r^*$.  Comparison of blocks now
gives the measure-and-prepare formula
\[
 F(X)=\sum_r \tr\bigl[(a_ra_r^*)^{\T}X\bigr]b_rb_r^*.
\]
Put $E_r=(a_ra_r^*)^{\T}\succeq0$.  For $W\succeq0$, each partial matrix
\[
 \tr_a\bigl[(I_k\otimes E_r)W\bigr]
 =\tr_a\bigl[(I_k\otimes\sqrt{E_r})W
                         (I_k\otimes\sqrt{E_r})\bigr]
\]
is positive.  Thus $(\mathrm{id}_k\otimes F)(W)$ is a sum of positive
products with second factors $b_rb_r^*$.  The converse follows by applying
the asserted property to $W=\Omega_a\Omega_a^*$.  This proves the
entanglement-breaking characterization \cite[Theorem 4]{HSR2003} at the
level of the CP cone.  The two transpose identities follow directly by
reindexing the matrix units in \eqref{eq:choi-definition}.
\end{proof}

\begin{lemma}\label{lem:product-range}
Every nonzero separable positive matrix has a nonzero product vector in its
range.  Consequently, a nonzero positive matrix whose range contains no
nonzero product vector is nonseparable.
\end{lemma}

\begin{proof}
Refine a separable decomposition by the spectral decompositions of its
positive factors.  It then has the form $Z=\sum_r v_rv_r^*$ with product
vectors $v_r=a_r\otimes b_r$.  Since
$\ker Z=\bigcap_r v_r^\perp$, every $v_r$ lies in
$(\ker Z)^\perp=\ran Z$.  If $Z\ne0$, at least one $v_r$ is nonzero.
\end{proof}

This is the range condition we use from the range criterion
\cite{Horodecki1997}.  It is only a necessary condition for separability;
the dimension-ten construction below establishes its failure directly.

\begin{definition}\label{def:represented-class}
A density operator $\omega$ on $\C^a\otimes\C^b$ belongs to the represented
class $\mathscr C$ if there are positive integers $r,s$, a vector
$\chi\in\C^r\otimes\C^s$, and linear maps
$\mathcal L:\Mat r\to\Mat a$ and $\mathcal M:\Mat s\to\Mat b$ such that
\begin{equation}\label{eq:represented-class}
 \omega=(\mathcal L\otimes\mathcal M)(\chi\chi^*),
 \qquad
 \mathcal L,\ \mathcal L\circ T_r,\ \mathcal M,\ \mathcal M\circ T_s
 \text{ are CP}.
\end{equation}
Only $\omega$ is required to have trace one.  The vector $\chi$ may be
unnormalized, and the maps need not preserve trace.
\end{definition}

This is a representation property, with no assumed physical implementation
of the representing maps.  Lemma~\ref{lem:ppt-stability} shows that the two
maps are PPT also in the output convention.  In particular the state itself
has positive partial transpose, since, for example,
$(T_a\otimes\mathrm{id}_b)(\omega)=
 ((T_a\circ\mathcal L)\otimes\mathcal M)(\chi\chi^*)\succeq0$.
The operational proof will use the representing maps themselves, in
addition to this partial-transpose consequence.

\begin{lemma}\label{lem:tensor-closure}
The class $\mathscr C$ is closed under finite tensor products, after
grouping all first-party systems and all second-party systems.  In
particular, $\omega^{\otimes n}\in\mathscr C$ for every
$\omega\in\mathscr C$ and every $n\geq1$.
\end{lemma}

\begin{proof}
Tensor the representing vectors and the first-party maps separately from
the second-party maps, and use the canonical local permutation unitaries to
group the systems.  Tensor products of CP maps are CP, and input transpose
in the grouped product basis is the tensor product of the individual input
transposes.  Hence the four maps required in
\eqref{eq:represented-class} remain CP.  The output is the tensor product of
the original states and has trace one.  The local permutation unitaries and
their conjugates are CP pre- and post-compositions, so they do not affect
the conclusion.
\end{proof}

\begin{lemma}[Choi transfer]\label{lem:choi-transfer}
Let $F:\Mat a\to\Mat b$ and $G:\Mat b\to\Mat c$ be CP maps, and set
$F^\sharp=T_a\circ F^\dagger\circ T_b$.  Then
\begin{equation}\label{eq:choi-transfer}
 J(G\circ F)=(F^\sharp\otimes G)(\Omega_b\Omega_b^*).
\end{equation}
If $F$ and $G$ are PPT and $Z=J(G\circ F)\ne0$, then
$\lambda:=\tr Z>0$ and
\begin{equation}\label{eq:choi-transfer-normalized}
 \frac{Z}{\lambda}
 =(F^\sharp\otimes G)
 \left(\frac{\Omega_b\Omega_b^*}{\lambda}\right)\in\mathscr C
 \quad\text{on }\C^a\otimes\C^c.
\end{equation}
\end{lemma}

\begin{proof}
Let $A_r\in\C^{b\times a}$ and $B_s\in\C^{c\times b}$ be Kraus matrices
for $F$ and $G$.  Transposing the adjoint formula gives
\[
 F^\sharp(Y)=\sum_r A_r^{\T}Y\overline{A_r}
            =\sum_r A_r^{\T}Y(A_r^{\T})^*.
\]
Thus the Kraus matrices of $F^\sharp:\Mat b\to\Mat a$ are the transposes
$A_r^{\T}$.  For each pair of Kraus matrices, direct expansion in the fixed
bases gives
\[
 (A_r^{\T}\otimes B_s)\Omega_b
       =(I_a\otimes B_sA_r)\Omega_a.
\]
Summing the corresponding rank-one matrices proves
\eqref{eq:choi-transfer}.  This is the rectangular CP form of the
transpose-adjoint identity in \cite[Lemma I.1(2)]{CMHW2019}; its use for
local products of PPT maps is discussed there in Section IV.B.

Suppose now that $F$ and $G$ are PPT.  Besides complete positivity of
$F^\sharp$ and $G$, Lemma~\ref{lem:ppt-stability} gives
\[
 F^\sharp\circ T_b=T_a\circ F^\dagger=(F\circ T_a)^\dagger
       \text{ CP},
 \qquad G\circ T_b\text{ CP}.
\]
These are exactly the input-transpose conditions for the two local maps in
Definition~\ref{def:represented-class}.  Finally,
\[
 Z=\sum_{r,s}|(I_a\otimes B_sA_r)\Omega_a\rangle
                 \langle(I_a\otimes B_sA_r)\Omega_a|,
 \qquad
 \lambda=\sum_{r,s}\norm{B_sA_r}_{\mathrm F}^{\,2}.
\]
The nonzero positive matrix $Z$ has $\lambda>0$.  Taking
$\chi=\Omega_b/\sqrt\lambda$ in \eqref{eq:represented-class} proves
\eqref{eq:choi-transfer-normalized}.  Its squared norm is $b/\lambda$,
which is allowed by the definition.  No trace-preserving hypothesis or
additional factor of $b$ is used.
\end{proof}
